Observe that
$\ang{1} = \ang{10} - \ang{9}$,
and $\ang{10} = \mfrac{\pi}{18}$.
Hence:
\begin{align*}
 \sin\paren{\ang{1}}
 &= \sin\paren{\ang{10} - \ang{9}}
  = \sin\paren{\ang{10}} \cos\paren{\ang{9}}
    - \cos\paren{\ang{10}} \sin\paren{\ang{9}} = \\
 &= \frac{1}{2 i}
    \paren{e^{i \frac{\pi}{18}} - e^{ -i \frac{\pi}{18}}}
    \cos\paren{\ang{9}}
    - \frac{1}{2} \paren{e^{i \frac{\pi}{18}} + e^{ -i \frac{\pi}{18}}}
    \sin\paren{\ang{9}} = \\
 &= \frac{1}{2}
    e^{i \frac{\pi}{18}}
    \paren{ - \sin\paren{\ang{9}} - i \cos\paren{\ang{9}}}
    + \frac{1}{2}
    e^{ -i \frac{\pi}{18}}
    \paren{ - \sin\paren{\ang{9}} + i \cos\paren{\ang{9}}}
\end{align*}
and:
\begin{align*}
 \cos\paren{\ang{1}}
 &= \cos\paren{\ang{10} - \ang{9}}
  = \cos\paren{\ang{10}} \cos\paren{\ang{9}}
    + \sin\paren{\ang{10}} \sin\paren{\ang{9}} = \\
 &= \frac{1}{2}
    \paren{e^{i \frac{\pi}{18}} + e^{ -i \frac{\pi}{18}}}
    \cos\paren{\ang{9}}
    + \frac{1}{2 i}
    \paren{e^{i \frac{\pi}{18}} - e^{ -i \frac{\pi}{18}}}
    \sin\paren{\ang{9}} = \\
 &= \frac{1}{2}
    e^{i \frac{\pi}{18}}
    \paren{ \cos\paren{\ang{9}} - i \sin\paren{\ang{9}}}
    + \frac{1}{2}
    e^{ -i \frac{\pi}{18}}
    \paren{\cos\paren{\ang{9}} + i \sin\paren{\ang{9}}}
\end{align*}
The term $e^{i \frac{\pi}{18}}$
is the principal cube root of
$e^{i \frac{\pi}{6}} = \mfrac{\sqrt{3}}{2} + \mfrac{1}{2} i$,
and the term $e^{- i \frac{\pi}{18}}$
is the principal cube root of
$e^{- i \frac{\pi}{6}} = \mfrac{\sqrt{3}}{2} - \mfrac{1}{2} i$,
so they can be written as:
\begin{alignat*}{2}
 & e^{i \frac{\pi}{18}}
 = \sqrt[3]{\frac{\sqrt{3} + i}{2}} \; , && \qquad
 e^{- i \frac{\pi}{18}}
 = \sqrt[3]{\frac{\sqrt{3} - i}{2}}
\end{alignat*}
The values of $\sin\paren{\ang{9}}$ and $\cos\paren{\ang{9}}$ are the following
(they can be calculated using the angle difference formulas
from $\ang{9} = \ang{45} - \ang{36}$):
\begin{align*}
 \sin\paren{\ang{9}}
 &= \frac{1}{8} \sqrt{2} \, \big(\sqrt{5} + 1 \big)
    - \frac{1}{4} \sqrt{5 - \sqrt{5}}
 \\
 \cos\paren{\ang{9}}
 &= \frac{1}{8} \sqrt{2} \, \big(\sqrt{5} + 1 \big)
    + \frac{1}{4} \sqrt{5 - \sqrt{5}}
\end{align*}
Altogether:
\begin{equation*}
\begin{split}
 \sin\paren{\ang{1}} \,
 &= \, \frac{1}{8} \, \sqrt[3]{\frac{\sqrt{3} + i}{2}}
   \, \cdot
   \Bigg( \!
     - \frac{1}{2} \sqrt{2} \, \big(\sqrt{5} + 1 \big)
     + \sqrt{5 - \sqrt{5}} \;\;
     - \,\; i \, \bigg(
     \frac{1}{2} \sqrt{2} \, \big(\sqrt{5} + 1 \big)
     + \sqrt{5 - \sqrt{5}} \,
     \bigg) \!
   \Bigg) \,\; +
 \\
   &\, + \, \frac{1}{8} \, \sqrt[3]{\frac{\sqrt{3} - i}{2}}
   \, \cdot
   \Bigg( \!
     - \frac{1}{2} \sqrt{2} \, \big(\sqrt{5} + 1 \big)
     + \sqrt{5 - \sqrt{5}} \;\;
     + \,\; i \, \bigg(
     \frac{1}{2} \sqrt{2} \, \big(\sqrt{5} + 1 \big)
     + \sqrt{5 - \sqrt{5}} \,
     \bigg) \!
   \Bigg)
\end{split}
\end{equation*}
and:
\begin{equation*}
\begin{split}
 \cos\paren{\ang{1}} \,
 &= \, \frac{1}{8} \, \sqrt[3]{\frac{\sqrt{3} + i}{2}}
   \, \cdot
   \Bigg( \,
     \frac{1}{2} \sqrt{2} \, \big(\sqrt{5} + 1 \big)
     + \sqrt{5 - \sqrt{5}} \;\;
     - \,\; i \, \bigg(
     \frac{1}{2} \sqrt{2} \, \big(\sqrt{5} + 1 \big)
     - \sqrt{5 - \sqrt{5}} \,
     \bigg) \!
   \Bigg) \,\; +
 \\
   &\, + \, \frac{1}{8} \, \sqrt[3]{\frac{\sqrt{3} - i}{2}}
   \, \cdot
   \Bigg( \,
     \frac{1}{2} \sqrt{2} \, \big(\sqrt{5} + 1 \big)
     + \sqrt{5 - \sqrt{5}} \;\;
     + \,\; i \, \bigg(
     \frac{1}{2} \sqrt{2} \, \big(\sqrt{5} + 1 \big)
     - \sqrt{5 - \sqrt{5}} \,
     \bigg) \!
   \Bigg)
\end{split}
\end{equation*}


For $\ang{2}$, observe that
$\ang{2} = \ang{20} - \ang{18}$,
and $\ang{20} = \mfrac{\pi}{9}$.
Hence:
\begin{align*}
 \sin\paren{\ang{2}}
 &= \sin\paren{\ang{20} - \ang{18}}
  = \sin\paren{\ang{20}} \cos\paren{\ang{18}}
    - \cos\paren{\ang{20}} \sin\paren{\ang{18}} = \\
 &= \frac{1}{2 i}
    \paren{e^{i \frac{\pi}{9}} - e^{ -i \frac{\pi}{9}}}
    \cos\paren{\ang{18}}
    - \frac{1}{2}
    \paren{e^{i \frac{\pi}{9}} + e^{ -i \frac{\pi}{9}}}
    \sin\paren{\ang{18}} = \\
 &= \frac{1}{2}
    e^{i \frac{\pi}{9}}
    \paren{ - \sin\paren{\ang{18}} - i \cos\paren{\ang{18}}}
    + \frac{1}{2}
    e^{ -i \frac{\pi}{9}}
    \paren{ - \sin\paren{\ang{18}} + i \cos\paren{\ang{18}}}
\end{align*}
and:
\begin{align*}
 \cos\paren{\ang{2}}
 &= \cos\paren{\ang{20} - \ang{18}}
  = \cos\paren{\ang{20}} \cos\paren{\ang{18}}
    + \sin\paren{\ang{20}} \sin\paren{\ang{18}} = \\
 &= \frac{1}{2}
    \paren{e^{i \frac{\pi}{9}} + e^{ -i \frac{\pi}{9}}}
    \cos\paren{\ang{18}}
    + \frac{1}{2 i}
    \paren{e^{i \frac{\pi}{9}} - e^{ -i \frac{\pi}{9}}}
    \sin\paren{\ang{18}} = \\
 &= \frac{1}{2}
    e^{i \frac{\pi}{9}}
    \paren{\cos\paren{\ang{18}} - i \sin\paren{\ang{18}}}
    + \frac{1}{2}
    e^{ -i \frac{\pi}{9}}
    \paren{\cos\paren{\ang{18}} + i \sin\paren{\ang{18}}}
\end{align*}
The term $e^{i \frac{\pi}{9}}$
is the principal cube root of
$e^{i \frac{\pi}{3}} = \mfrac{1}{2} + i \mfrac{\sqrt{3}}{2}$,
and the term $e^{- i \frac{\pi}{9}}$
is the principal cube root of
$e^{- i \frac{\pi}{3}} = \mfrac{1}{2} - i \mfrac{\sqrt{3}}{2}$,
so they can be written as:
\begin{alignat*}{2}
 & e^{i \frac{\pi}{9}}
 = \sqrt[3]{\frac{1 + i \sqrt{3}}{2}} \; , && \qquad
 e^{- i \frac{\pi}{9}}
 = \sqrt[3]{\frac{1 - i \sqrt{3}}{2}}
\end{alignat*}
The values of $\sin\paren{\ang{18}}$ and $\cos\paren{\ang{18}}$ are the following:
\begin{align*}
 \sin\paren{\ang{18}}
 &= \frac{1}{8} \big(\sqrt{5} - 1 \big)
 \\
 \cos\paren{\ang{18}}
 &= \frac{1}{4} \sqrt{2} \, \sqrt{5 + \sqrt{5}}
\end{align*}
Altogether:
\begin{equation*}
 \sin\paren{\ang{2}}
 \, = \, \frac{1}{8} \, \sqrt[3]{\frac{1 + i \sqrt{3}}{2}}
   \, \cdot
   \Bigg( \!\!
     - \frac{\sqrt{5} - 1}{2}
     \, - \, i \, \sqrt{2} \, \sqrt{5 + \sqrt{5}} \,
   \Bigg)
 \, + \;
   \frac{1}{8} \, \sqrt[3]{\frac{1 - i \sqrt{3}}{2}}
   \, \cdot
   \Bigg( \!\!
     - \frac{\sqrt{5} - 1}{2}
     \, + \, i \, \sqrt{2} \, \sqrt{5 + \sqrt{5}} \,
   \Bigg)
\end{equation*}
and:
\begin{equation*}
 \cos\paren{\ang{2}}
 \, = \, \frac{1}{8} \, \sqrt[3]{\frac{1 + i \sqrt{3}}{2}}
   \, \cdot
   \Bigg( \!
     \sqrt{2} \, \sqrt{5 + \sqrt{5}}
     \, - \, i \, \frac{\sqrt{5} - 1}{2}
   \Bigg)
 \, + \;
   \frac{1}{8} \, \sqrt[3]{\frac{1 - i \sqrt{3}}{2}}
   \, \cdot
   \Bigg( \!
     \sqrt{2} \, \sqrt{5 + \sqrt{5}}
     \, + \, i \, \frac{\sqrt{5} - 1}{2}
   \Bigg)
\end{equation*}
